\textbf{Phần 1. Trắc nghiệm lựa chọn}
\setcounterpageref{ex}{0}

\begin{ex} %[1D7V26]
Xăng-ti-mét viết tắt là chữ gì trong bảng chữ cái tiếng Anh? Số chữ cái viết tắt là bao nhiêu?
\choice
{$1\text{ cm}$}
{\True $2\text{ cm}$}
{$3\text{ cm}$}
{$4\text{ cm}$}
\loigiai{
Xăng-ti-mét được viết tắt là cm (gồm 2 chữ cái c và m).
\textbf{Đáp số:} $2$.\par
\textbf{Đáp án của câu hỏi là B}
}
\end{ex}

\begin{ex} %[1D7V26]
Tính: $8\text{ cm} + 2\text{ cm}$. Kết quả là bao nhiêu xăng-ti-mét?
\choice
{$0\text{ cm}$}
{$20\text{ cm}$}
{$9\text{ cm}$}
{\True $10\text{ cm}$}
\loigiai{
Ta có: $8\text{ cm} + 2\text{ cm} = 10\text{ cm}$.
\textbf{Đáp số:} $10$.\par
\textbf{Đáp án của câu hỏi là D}
}
\end{ex}

\begin{ex} %[1D7V26]
Tính: $10\text{ cm} - 3\text{ cm}$. Kết quả là bao nhiêu xăng-ti-mét?
\choice
{\True $7\text{ cm}$}
{$6\text{ cm}$}
{$8\text{ cm}$}
{$5\text{ cm}$}
\loigiai{
Ta có: $10\text{ cm} - 3\text{ cm} = 7\text{ cm}$.
\textbf{Đáp số:} $7$.\par
\textbf{Đáp án của câu hỏi là A}
}
\end{ex}

\begin{ex} %[1D7V26]
Tính: $12\text{ cm} + 5\text{ cm}$. Kết quả là bao nhiêu xăng-ti-mét?
\choice
{$16\text{ cm}$}
{$18\text{ cm}$}
{\True $17\text{ cm}$}
{$15\text{ cm}$}
\loigiai{
Ta có: $12\text{ cm} + 5\text{ cm} = 17\text{ cm}$.
\textbf{Đáp số:} $17$.\par
\textbf{Đáp án của câu hỏi là C}
}
\end{ex}

\begin{ex} %[1D7V26]
Tính: $19\text{ cm} - 7\text{ cm}$. Kết quả là bao nhiêu xăng-ti-mét?
\choice
{$11\text{ cm}$}
{\True $12\text{ cm}$}
{$13\text{ cm}$}
{$10\text{ cm}$}
\loigiai{
Ta có: $19\text{ cm} - 7\text{ cm} = 12\text{ cm}$.
\textbf{Đáp số:} $12$.\par
\textbf{Đáp án của câu hỏi là B}
}
\end{ex}

\begin{ex} %[1D7V26]
Điền số thích hợp: $5\text{ cm} + \dots\text{ cm} = 9\text{ cm}$.
\choice
{$3\text{ cm}$}
{$5\text{ cm}$}
{$2\text{ cm}$}
{\True $4\text{ cm}$}
\loigiai{
Số cần điền là: $9 - 5 = 4$.
\textbf{Đáp số:} $4$.\par
\textbf{Đáp án của câu hỏi là D}
}
\end{ex}

\begin{ex} %[1D7V26]
Điền số thích hợp: $15\text{ cm} - \dots\text{ cm} = 10\text{ cm}$.
\choice
{\True $5\text{ cm}$}
{$4\text{ cm}$}
{$6\text{ cm}$}
{$3\text{ cm}$}
\loigiai{
Số cần điền là: $15 - 10 = 5$.
\textbf{Đáp số:} $5$.\par
\textbf{Đáp án của câu hỏi là A}
}
\end{ex}

\begin{ex} %[1D7V26]
Một chiếc bút sáp màu dài 8 cm, chiếc bút lông dài 14 cm. Hỏi chiếc bút lông dài hơn chiếc bút sáp bao nhiêu xăng-ti-mét?
\choice
{$5\text{ cm}$}
{$7\text{ cm}$}
{\True $6\text{ cm}$}
{$4\text{ cm}$}
\loigiai{
Bút lông dài hơn bút sáp là: $14 - 8 = 6$ (cm).
\textbf{Đáp số:} $6$.\par
\textbf{Đáp án của câu hỏi là C}
}
\end{ex}

\begin{ex} %[1D7V26]
Đoạn thẳng AB dài 6 cm, đoạn thẳng BC dài 4 cm. Ba điểm A, B, C thẳng hàng theo thứ tự đó. Hỏi đoạn thẳng AC dài bao nhiêu xăng-ti-mét?
\choice
{$0\text{ cm}$}
{\True $10\text{ cm}$}
{$20\text{ cm}$}
{$9\text{ cm}$}
\loigiai{
Độ dài đoạn thẳng AC là: $6 + 4 = 10$ (cm).
\textbf{Đáp số:} $10$.\par
\textbf{Đáp án của câu hỏi là B}
}
\end{ex}

\begin{ex} %[1D7V26]
Một sợi dây đồng dài 18 cm. Bạn Hùng cắt lấy 8 cm để làm thí nghiệm. Hỏi sợi dây đồng còn lại dài bao nhiêu xăng-ti-mét?
\choice
{$0\text{ cm}$}
{$20\text{ cm}$}
{$9\text{ cm}$}
{\True $10\text{ cm}$}
\loigiai{
Độ dài sợi dây đồng còn lại là: $18 - 8 = 10$ (cm).
\textbf{Đáp số:} $10$.\par
\textbf{Đáp án của câu hỏi là D}
}
\end{ex}

\begin{ex} %[1D7V26]
Tính: $7\text{ cm} + 3\text{ cm} + 4\text{ cm}$. Kết quả là bao nhiêu xăng-ti-mét?
\choice
{\True $14\text{ cm}$}
{$13\text{ cm}$}
{$15\text{ cm}$}
{$12\text{ cm}$}
\loigiai{
Thực hiện: $7 + 3 = 10$; $10 + 4 = 14$ (cm).
\textbf{Đáp số:} $14$.\par
\textbf{Đáp án của câu hỏi là A}
}
\end{ex}

\begin{ex} %[1D7V26]
Tính: $18\text{ cm} - 5\text{ cm} - 3\text{ cm}$. Kết quả là bao nhiêu xăng-ti-mét?
\choice
{$0\text{ cm}$}
{$20\text{ cm}$}
{\True $10\text{ cm}$}
{$9\text{ cm}$}
\loigiai{
Thực hiện: $18 - 5 = 13$; $13 - 3 = 10$ (cm).
\textbf{Đáp số:} $10$.\par
\textbf{Đáp án của câu hỏi là C}
}
\end{ex}

\textbf{Phần 2. Trắc nghiệm đúng sai}
\setcounterpageref{ex}{0}

\begin{ex} %[1D7V26]
Xét tính đúng hoặc sai của các khẳng định toán học sau đây:
\choiceTF
{Một cây thước dài 20 cm bị gãy mất một đầu dài 4 cm. đoạn thước còn lại dài $17\text{ cm}$ xăng-ti-mét.}
{\True Con sâu đo bò được 5 cm trên chiếc lá, rồi bò thêm 6 cm nữa đến cuống lá. con sâu đo đã bò được tất cả $11\text{ cm}$ xăng-ti-mét.}
{Một hộp bút dài 17 cm, một cục tẩy dài 3 cm. cần $15\text{ cm}$ xăng-ti-mét nữa để cục tẩy dài bằng hộp bút.}
{\True Đoạn thẳng MN dài 13 cm, đoạn thẳng PQ ngắn hơn đoạn MN là 3 cm. đoạn thẳng PQ dài $10\text{ cm}$ xăng-ti-mét.}
\loigiai{
\begin{itemize}
    \item Ý a) Sai: Kết quả đúng là $16$, không phải $17$. Đoạn thước còn lại dài: $20 - 4 = 16$ (cm).
\textbf{Đáp số:} $16$.
    \item Ý b) Đúng: Tổng quãng đường sâu đo bò là: $5 + 6 = 11$ (cm).
\textbf{Đáp số:} $11$.
    \item Ý c) Sai: Kết quả đúng là $14$, không phải $15$. Hiệu độ dài là: $17 - 3 = 14$ (cm).
\textbf{Đáp số:} $14$.
    \item Ý d) Đúng: Độ dài đoạn thẳng PQ là: $13 - 3 = 10$ (cm).
\textbf{Đáp số:} $10$.
\end{itemize}
\begin{itemize}
    \item a) \textbf{Sai}.
    \item b) \textbf{Đúng}.
    \item c) \textbf{Sai}.
    \item d) \textbf{Đúng}.
\end{itemize}
}
\end{ex}

\begin{ex} %[1D7V26]
Xét tính đúng hoặc sai của các bài toán và phép tính sau đây:
\choiceTF
{\True Kết quả của bài toán «Tổng độ dài của hai đoạn thẳng MN và PQ ở bài toán trên là bao nhiêu xăng-ti-mét» là $23\text{ cm}$.}
{\True Một băng giấy dài 15 cm. Bạn Mai cắt băng giấy thành hai đoạn, trong đó có một đoạn dài 8 cm. đoạn còn lại dài $7\text{ cm}$ xăng-ti-mét.}
{Giá trị của phép tính $20\text{ cm} - 10\text{ cm} + 5\text{ cm}$ là $14\text{ cm}$.}
{Một hộp bút dài 17 cm, một cục tẩy dài 3 cm. cần $15\text{ cm}$ xăng-ti-mét nữa để cục tẩy dài bằng hộp bút.}
\loigiai{
\begin{itemize}
    \item Ý a) Đúng: Tổng độ dài là: $13 + 10 = 23$ (cm).
\textbf{Đáp số:} $23$.
    \item Ý b) Đúng: Đoạn băng giấy còn lại dài: $15 - 8 = 7$ (cm).
\textbf{Đáp số:} $7$.
    \item Ý c) Sai: Kết quả đúng là $15$, không phải $14$. Thực hiện: $20 - 10 = 10$; $10 + 5 = 15$ (cm).
\textbf{Đáp số:} $15$.
    \item Ý d) Sai: Kết quả đúng là $14$, không phải $15$. Hiệu độ dài là: $17 - 3 = 14$ (cm).
\textbf{Đáp số:} $14$.
\end{itemize}
\begin{itemize}
    \item a) \textbf{Đúng}.
    \item b) \textbf{Đúng}.
    \item c) \textbf{Sai}.
    \item d) \textbf{Sai}.
\end{itemize}
}
\end{ex}

\textbf{Phần 3. Trả lời ngắn}
\setcounterpageref{ex}{0}

\begin{ex} %[1D7V26]
Một con cào cào nhảy bước thứ nhất được 8 cm, bước thứ hai nhảy được 9 cm. Hỏi cả hai bước cào cào nhảy được bao nhiêu xăng-ti-mét?
\shortans{17}
\loigiai{
\textbf{Đáp số:} $17$\par
Cả hai bước cào cào nhảy được: $8 + 9 = 17$ (cm).
\textbf{Đáp số:} $17$.
}
\end{ex}

\begin{ex} %[1D7V26]
Chiều dài chiếc điện thoại di động là 14 cm, chiều rộng là 7 cm. Chiều dài hơn chiều rộng bao nhiêu xăng-ti-mét?
\shortans{7}
\loigiai{
\textbf{Đáp số:} $7$\par
Chiều dài hơn chiều rộng là: $14 - 7 = 7$ (cm).
\textbf{Đáp số:} $7$.
}
\end{ex}

\begin{ex} %[1D7V26]
Tìm số thích hợp: $11\text{ cm} + \dots\text{ cm} = 16\text{ cm}$.
\shortans{5}
\loigiai{
\textbf{Đáp số:} $5$\par
Số cần điền là: $16 - 11 = 5$.
\textbf{Đáp số:} $5$.
}
\end{ex}

\begin{ex} %[1D7V26]
Tìm số thích hợp: $19\text{ cm} - \dots\text{ cm} = 11\text{ cm}$.
\shortans{8}
\loigiai{
\textbf{Đáp số:} $8$\par
Số cần điền là: $19 - 11 = 8$.
\textbf{Đáp số:} $8$.
}
\end{ex}

\begin{ex} %[1D7V26]
Một cây nến cao 16 cm. Sau một giờ cháy bớt đi 6 cm. Hỏi cây nến còn lại cao bao nhiêu xăng-ti-mét?
\shortans{10}
\loigiai{
\textbf{Đáp số:} $10$\par
Chiều cao còn lại của cây nến là: $16 - 6 = 10$ (cm).
\textbf{Đáp số:} $10$.
}
\end{ex}

\begin{ex} %[1D7V26]
Đoạn dây thép thứ nhất dài 12 cm, đoạn dây thép thứ hai dài 7 cm. Nối hai đoạn lại (không tính mối nối), sợi dây mới dài bao nhiêu xăng-ti-mét?
\shortans{19}
\loigiai{
\textbf{Đáp số:} $19$\par
Chiều dài sợi dây mới là: $12 + 7 = 19$ (cm).
\textbf{Đáp số:} $19$.
}
\end{ex}
